\documentclass[10pt]{article}

\usepackage[top=0.5in,bottom=0.5in,left=0.7in,right=0.7in]{geometry}
\usepackage[T1]{fontenc}
\usepackage{lmodern}
\usepackage{microtype}

\title{SteerEats: Explaining and Steering Restaurant Recommendations in NYC\vspace{-.5em}}
\author{Fall 2026, DS-GA 3001 Project Update\vspace{-.5em}}
\date{Team 4: Fangyu Qu (fq2069@nyu.edu), Yila Cao (yc8155@nyu.edu)}

\newcommand{\placeholder}[1]{\textbf{[TO FILL: #1]}}
\setlength{\parindent}{0pt}
\setlength{\parskip}{1pt}
\setlength{\leftmargini}{1.2em}

% Keep standard research-paper headings while using compact spacing for one page.
\makeatletter
\renewcommand\section{\@startsection{section}{1}{\z@}%
  {-0.95ex \@plus -.2ex \@minus -.1ex}%
  {0.45ex \@plus .1ex}%
  {\normalfont\large\bfseries}}
\makeatother

\begin{document}
\maketitle
\thispagestyle{empty}
\vspace{-2.5em}

\section{Progress Since the Proposal}
We narrowed the project to an \textbf{individual recommender}. Its central visualization claim is that a linked
before-and-after view shows why a ranking changed in ways that aggregate ranking metrics cannot. We dropped ALS,
the UMAP and calibration views and the separate researcher dashboard; group recommendation and steering cost
are now stretch goals.

\textbf{Data.} In the Google Local New York data (33.5M reviews), 24,857 restaurants lie inside NYC's 2020
neighborhoods, with 6.9M reviews by 2.2M users up to September 2021. Our
development sample is 20,000 random users (seed 42) with at least 10 NYC restaurant reviews since 2017: 422,788
first reviews of 21,155 restaurants. We split it by time: validation from December 2019, test from August 2020
(the last 10\% of interactions).

\textbf{Matching.} 42.4\% of NYC restaurants (51.5\% of those open at the 2021 crawl) match a DOHMH inspection record
by name, house number and distance; they receive 62\% of all reviews. In a hand check of 100 random matches, 97 were
correct, 1 wrong and 2 unclear. Most misses are closures: the inspection file lists only recently active restaurants, and
for half of the unmatched open restaurants another business now uses the address.

\textbf{Baselines.} On the test period (8,449 users), item-based CF reaches Recall@10 $=$ 0.026 and NDCG@10 $=$ 0.018,
about twice global popularity (0.012 and 0.008); popularity over the last 12 months beats all-time popularity
(Recall@10 $=$ 0.014). \placeholder{first logistic-regression result, or delete this sentence}

\textbf{Ratings over time.} Google's displayed rating is the all-time mean of the reviews (within 0.2 stars for all
17,614 restaurants with at least 50 reviews), so it hides change: for 5.5\% of the 8,311 restaurants with at least 20
reviews in the last year of data, the recent mean is at least 0.5 stars away from the all-time mean (225 down, 229 up).

\section{Key Paper and Gap}
TasteWeights (Bostandjiev et al., RecSys 2012) shows a user's liked artists, three sources and the recommendations as
linked columns, with sliders on items and sources; its user study found that control and explanation raised satisfaction. We keep its linked input--model--output view and direct control. But it explains the current list, not the
change: users cannot see which items entered, left or moved after steering; its model is a weighted sum, not a
learned ranker; and it ignores location.

\section{Response to Feedback}
\textbf{What a rating means.} A star rating reflects expectations about both quality and cost, so the interface shows
price level next to each rating and compares restaurants within a price level, along with review count and a
rolling mean with its uncertainty over time.

\textbf{Time as a control.} Instead of a separate drift study, users choose a date range (e.g., the last 12
months); restaurant features are recomputed from reviews in that range, and the before-and-after view explains the
resulting rank changes.

\textbf{Explanation accounting.} For each item, the interface applies the stages in a fixed order (learned model,
distance and popularity adjustments, diversity re-ranking), compares the reconstructed rank change with the observed
change and reports the residual. SHAP explains the raw model score, not the final rank.

\textbf{Evidence.} We use worked case studies and a small expert or formative review instead of a controlled study that needs
IRB approval. LLM summaries of review themes are optional context. Bursts of extreme ratings are labeled
only as anomalies, not as fake reviews.

\section{Scope, Plan and Division of Labor}
The core is item-based CF and popularity baselines, then a learned ranker (logistic regression, then LightGBM if
feasible), evaluated with Recall@10 and NDCG@10 across time windows. The linked interface has a map, before-and-after
ranks, score and rank decomposition, rating history with a date-range control, and relevance--distance--diversity
trade-offs.

\textbf{Fangyu} owns data preparation (NYC filtering, time split, inspection matching, data dictionary), the baselines,
and the rating-history and coverage views. \textbf{Yila} owns the learned ranker and explanation pipeline and the
rank-change, SHAP, control, diversity and map views. Both review code, integrate the interface, run case
studies and write the report.

\textbf{Timeline.} Nov 3--9: learned ranker and rank-change decomposition; Nov 10--16: linked interface;
Nov 17--23: case studies and expert review; Nov 24--30: slides and demo; Dec 1--14: final report.
\textbf{GitHub:} \texttt{https://github.com/fangyu-qu/SteerEats}

\end{document}
